\chapter{Iteration 7 --- what can EEG features decode about the other modalities?}
\label{ch:iter7}

\section{Task and inputs}

\paragraph{Question.} If we take the same iter-5 368-dimensional EEG
spectral feature vector that decodes attended-speaker hemisphere at
$\sim 0.72$, and instead ask it to predict \emph{each individual
feature} of gaze / IMU / video / audio, how well does it do?

This is a direct readout of which external-modality features are
linearly regressable from EEG spectral content. High $r$ for a feature
means EEG carries that feature's variance in a linearly accessible way
--- i.e.\ the EEG is ``about'' that dimension.

\paragraph{Method.} Per subject, per modality, per target feature:
\begin{enumerate}
  \item Align per-trial feature tables on \code{(subject, trial)} ---
        audio features are trial-level, so they're broadcast across
        subjects.
  \item $X$ = 368-D EEG spectral features (standardised), $y$ = one
        scalar target feature.
  \item 5-fold KFold ridge regression ($\lambda = 10$), record
        held-out Pearson $r$ between predicted and actual $y$.
  \item Pool across 16 subjects with bootstrap 95\,\% CI.
\end{enumerate}

\paragraph{Target feature counts.}
23 gaze / 35 IMU / 15 video / 20 audio = 93 targets; per-subject-per-target
CV gives $16 \times 93 = 1488$ ridge fits, all on CPU in under a minute.

\section{Top EEG-decodable features per modality}

\subsection{Gaze (23 features available)}

\begin{table}[ht]
\centering
\begin{tabular}{lrr}
\toprule
Gaze feature                  & Pooled $r$ & 95\,\% CI \\
\midrule
\code{sacc\_rate}             & \textbf{+0.549} & [+0.467, +0.630] \\
\code{gx\_valid}              & +0.540         & [+0.404, +0.672] \\
\code{R\_valid}               & +0.508         & [+0.385, +0.632] \\
\code{L\_valid}               & +0.472         & [+0.314, +0.621] \\
\code{R\_pup}                 & +0.419         & [+0.323, +0.512] \\
\code{L\_pup}                 & +0.385         & [+0.261, +0.500] \\
\code{L\_az\_std}             & +0.361         & [+0.267, +0.467] \\
\code{R\_el\_mean}            & +0.342         & [+0.262, +0.427] \\
\code{gx\_std}                & +0.334         & [+0.232, +0.449] \\
\code{R\_az\_mean}            & +0.332         & [+0.208, +0.460] \\
\code{gy\_mean}               & +0.311         & [+0.221, +0.409] \\
\bottomrule
\end{tabular}
\end{table}

\noindent
The strongest EEG$\to$gaze predictor is \textbf{saccade rate}
($r=+0.549$). Per-eye validity (fraction of samples where the
eye-tracker locked a gaze vector) and pupil diameter are also highly
linearly predictable --- these reflect arousal / eye-opening, which
the EEG correlates with through both alpha and frontal power.

\subsection{IMU (35 features available)}

\begin{table}[ht]
\centering
\begin{tabular}{lrr}
\toprule
IMU feature                   & Pooled $r$ & 95\,\% CI \\
\midrule
\code{gyr\_mag\_mean}         & \textbf{+0.559} & [+0.468, +0.647] \\
\code{total\_rotation}        & +0.556         & [+0.465, +0.644] \\
\code{gyr\_mag\_p95}          & +0.541         & [+0.445, +0.634] \\
\code{gyr\_mag\_std}          & +0.534         & [+0.444, +0.621] \\
\code{gx\_std}                & +0.529         & [+0.443, +0.616] \\
\code{az\_mean}               & +0.507         & [+0.412, +0.609] \\
\code{az\_peak}               & +0.506         & [+0.401, +0.602] \\
\code{az\_std}                & +0.492         & [+0.422, +0.562] \\
\code{gy\_std}                & +0.466         & [+0.362, +0.571] \\
\code{acc\_mag\_std}          & +0.461         & [+0.376, +0.552] \\
\bottomrule
\end{tabular}
\end{table}

\noindent
\textbf{Head-motion amplitude features are the most linearly predictable
from EEG} — integrated rotation, gyroscope magnitude mean / std / 95th
pct, all above $r = 0.5$. This is not surprising: EMG contamination in
temporal / frontotemporal EEG channels is a direct sensor readout of
neck and head muscle activity.

\subsection{Video (15 features available)}

\begin{table}[ht]
\centering
\begin{tabular}{lrr}
\toprule
Video feature                 & Pooled $r$ & 95\,\% CI \\
\midrule
\code{motion\_energy\_mean}   & \textbf{+0.487} & [+0.403, +0.565] \\
\code{motion\_energy\_std}    & +0.471         & [+0.389, +0.550] \\
\code{motion\_energy\_peak}   & +0.455         & [+0.387, +0.527] \\
\code{brightness\_std}        & +0.453         & [+0.363, +0.545] \\
\code{flow\_mag\_mean}        & +0.452         & [+0.347, +0.554] \\
\code{flow\_std\_mean}        & +0.424         & [+0.327, +0.521] \\
\code{brightness\_mean}       & +0.352         & [+0.267, +0.431] \\
\code{flow\_mag\_std}         & +0.329         & [+0.232, +0.422] \\
\code{flow\_mag\_peak}        & +0.277         & [+0.193, +0.363] \\
\code{flow\_dir\_consistency} & +0.063         & [$-0.031$, +0.158] \\
\bottomrule
\end{tabular}
\end{table}

\noindent
Scene motion-energy and optical-flow magnitudes are linearly decodable
from EEG at $r \approx 0.45$. Directional flow features (\code{flow\_dir\_consistency},
\code{flow\_dir\_cy}) are at zero --- EEG doesn't know which way the
camera panned, only how much. This matches the IMU story: EEG encodes
motion amplitude, not motion direction.

\subsection{Audio (20 features available)}

\begin{table}[ht]
\centering
\begin{tabular}{lrr}
\toprule
Audio feature                 & Pooled $r$ & 95\,\% CI \\
\midrule
\code{una\_env\_der\_pos}     & $-0.086$ & [$-0.139$, $-0.038$] \\
\code{una\_env\_std}          & $-0.065$ & [$-0.121$, $-0.011$] \\
\code{att\_zcr}               & $+0.057$ & [$+0.002$, $+0.116$] \\
\code{una\_centroid}          & $+0.051$ & [$-0.009$, $+0.107$] \\
\code{una\_zcr}               & $+0.038$ & [$-0.012$, $+0.092$] \\
\code{att\_env\_mean}         & $-0.036$ & [$-0.071$, $+0.001$] \\
\code{att\_env\_std}          & $+0.031$ & [$-0.032$, $+0.092$] \\
\code{una\_flatness}          & $+0.027$ & [$-0.031$, $+0.082$] \\
\code{una\_rms}               & $-0.020$ & [$-0.080$, $+0.034$] \\
\code{att\_centroid}          & $+0.020$ & [$-0.026$, $+0.068$] \\
\bottomrule
\end{tabular}
\end{table}

\noindent
\textbf{EEG does not linearly predict any summary audio feature}. Every
correlation is within $\pm 0.09$, and the few that cross zero do so
weakly. This is neither surprising nor damning --- attended-speech
envelope \emph{tracking} is a time-resolved operation (the TRF works on
aligned time-series, not on per-trial scalar summaries), and the iter-5
spectral features deliberately summarise over the 30-s trial, collapsing
any trial-internal tracking to zero-mean. What this result does confirm
is that per-trial audio summary statistics are \emph{not} the axis along
which EEG spectral content varies --- the AAD signal is not
``EEG looks like the spectrogram''.

\section{Joint picture}

\begin{center}
\begin{tabular}{lcc}
\toprule
Modality & Top feature & Top $r$ \\
\midrule
Gaze     & \code{sacc\_rate}          & +0.549 \\
IMU      & \code{gyr\_mag\_mean}      & +0.559 \\
Video    & \code{motion\_energy\_mean} & +0.487 \\
Audio    & \code{una\_env\_der\_pos}  & $-0.086$ \\
\bottomrule
\end{tabular}
\end{center}

EEG spectral content is \textbf{strongly linearly related to head
motion, gaze validity / amplitude, and scene-motion energy}, and
\textbf{essentially unrelated to per-trial audio summary statistics}.

\paragraph{Reconciling with iter-5 and iter-6.}
This looks at first like a contradiction --- if EEG shares $r \approx
0.5$ with IMU and gaze, isn't the AAD signal just motion in disguise?
But iter-6 (§\ref{ch:iter6}) already showed that regressing gaze + IMU
out of EEG drops hemisphere AAD by only $-0.7$\,pp. The two findings
together tell a clean story:

\begin{itemize}
  \item EEG has a substantial component that \emph{is} head-motion /
        EMG artefact — this is real and predictable from external sensors.
  \item The \emph{attended-speaker-discriminating} component of EEG is
        a smaller but linearly separable signal that lives in
        central/parietal channels and spectral/connectivity features
        (iter-5 + iter-6). It survives artefact regression because the
        regressor subspace and the attention subspace overlap only
        partially.
\end{itemize}

\paragraph{Implications for iter-8.}
\begin{enumerate}
  \item Build an \textbf{artefact-subtraction decoder}: first predict
        gaze / IMU / video from EEG (the trivial $r \approx 0.5$
        regression), then train AAD on the \emph{residuals} of EEG
        after subtracting what external sensors could have predicted.
        This is a subject-level orthogonalisation that should isolate
        the audio-cortex contribution more cleanly than channel-wise
        regression did in iter-6.
  \item The zero predictability of audio features from EEG spectral
        summaries reinforces that we should not expect per-trial
        scalar summaries of audio to help an EEG decoder --- the
        temporal-tracking axis (TRF / CCA on time-series) is the
        right place to find the EEG-audio link, despite its weak
        pooled numbers in iter-3 and iter-4.
\end{enumerate}
